%!TEX root = ../main.tex

\section{Введение}

Электрический пробой -- это явление резкого возрастания тока в диэлектрике при приложении электрического напряжения выше некоторого критического значения. Механизм разрушения диэлектрика под действием электрического поля сложен и многообразен: оно может иметь различные причины, характер развития, сопутствующие физические процессы \cite{vorobiev_dielectric_physics}. В настоящей работе рассматривается случай твердого диэлектрика.

В силу сложности явления не удается найти математическую модель, чтобы описать его непосредственно, поэтому делается попытка построить приближенную модель, адекватно отражающую явление. Для этого научным коллективом был избран метод диффузной границы; модель развития канала электрического пробоя типа диффузной границы предложена в статье \cite{pitike_dielectric_breakdown}.

Модели типа диффузной границы в настоящее время составляют целый класс подходов для решения прикладных задач гидродинамики \cite{lamorgese_flow_modeling, kim_fluid_flows, xu_hydrodynamics}, механики деформируемого тела и теории трещин \cite{ambati_fracture}, материаловедения \cite{provatas_materials}, солидификации и теории фазовых переходов \cite{boettinger_solidification, cartalade_phase_separation, gransaly_solidification}, описания кристаллических структур \cite{emmerich_crystal, asadi_crystal, provatas_crystal}. В частности, предложенная в статье \cite{pitike_dielectric_breakdown} модель получена обобщением подобной модели распространения плоской трещины в упругой среде.

Упомянутая модель исследовалась научным коллективом с разных сторон, что изложено в ряде статей. В оригинальной работе \cite{pitike_dielectric_breakdown} модель сформулирована для двух первичных переменных -- фазового поля $\phi$ и электрического потенциала $\Phi$. В статье \cite{zipunova_computational} рассматривается также третья переменная -- пространственная плотность $\rho$ электрического заряда; именно в такой форме модель будет исследована в настоящей работе. Статья \cite{zipunova_conservative} посвящена еще более развитой модели с учетом температуры $\theta$, статья \cite{zipunova_thermomechanical} -- к тому же с учетом плотности и скорости подвижной среды. В работе \cite{ponomarev_stability} произведен анализ устойчивости положений равновесия модели в упрощенной, одномерной, постановке и дано содержательное условие устойчивости явной разностной схемы для задачи. В работах \cite{zipunova_higher_codimension, ponomarev_generalization} исследуется корректность описания фазовым полем одномерного объекта канала пробоя в трехмерном пространстве; предлагается обобщение модели для существования решений у ряда модельных краевых задач.

Кратко перечислим основные положения метода диффузной границы в применении к задаче развития канала электрического пробоя. Еще раз отметим, что модель феноменологическая: без рассмотрения причин образования пробоя на микроуровне строятся уравнения, описывающие эволюцию состояния среды, согласующиеся с базовыми законами физики. Вещество в системе находится в разных состояниях -- фазах, -- то есть не однородно, причем части вещества в одном состоянии образуют некоторые однородные области. С точки зрения физики <<раздел>> фаз имеет пренебрежимо малую толщину, его можно считать плоской поверхностью. Однако в моделях типа диффузной границы принимается, что распределение фаз в пространстве задано гладкой функцией~$\phi$, называемой фазовым полем: внутри каждой области однородности значение $\phi$ близко к постоянному, а на разделяющем слое (<<диффузной границе>>) меняется пусть и быстро, но непрерывно. В рассматриваемой модели вещество имеет две различные фазы: <<неповрежденное>> -- $\phi \approx 1$ -- и <<полностью разрушенное>> (то есть относящееся к каналу пробоя) -- $\phi \approx 0$. Характерная толщина раздела фаз определяется параметрами модели. На разрушение среды, то есть на постепенный переход вещества от состояния $\phi \approx 1$ к состоянию $\phi \approx 0$, тратится энергия электрического поля.

\absent{Продолжение}